% Source: https://artofproblemsolving.com/community/c3h1199872_a_rightangle_triangle
% Stable identity: aops-topic-1199872-q01
Triangle $APQ$ is similar to triangle $ABC$ with $P\leftrightarrow B$ and $Q\leftrightarrow C$, so for some $t>0$, $AQ=tb$, $AP=tc$ and $PQ=ta$, where $a=\sqrt{b^2+c^2}$. Put $A=(0,0)$, $B=(c,0)$, $C=(0,b)$, so $Q=(tb,0)$, $P=(0,tc)$, and $BC$ is $bx+cy=bc$.

\begin{asy}
usepackage("amsmath");
size(4cm);
defaultpen(fontsize(10pt));
pair A=(0,0),B=(4,0),C=(0,3); real t=0.45; pair Q=(t*3,0),P=(0,t*4); pair Ee=P+dot(B-P,unit(C-B))*unit(C-B)*0;
pair u=unit(C-B); pair n=(u.y,-u.x); pair E2=P-dot(P-B,n)*n, F2=Q-dot(Q-B,n)*n;
draw(A--B--C--cycle,linewidth(0.9)); draw(P--Q,blue+linewidth(0.9)); draw(P--E2,dashed); draw(Q--F2,dashed);
label("$A$",A,SW);label("$B$",B,SE);label("$C$",C,N);label("$P$",P,W);label("$Q$",Q,S);label("$E$",E2,NE);label("$F$",F2,NE);
label("\mbox{}",(min(currentpicture,user=true)+max(currentpicture,user=true))/2,black); // keeps the text after the figure black when LaTeX typesets the labels inline
\end{asy}

The distances to $BC$ are
\[PE=\dfrac{bc-c\cdot tc}{a},\qquad QF=\dfrac{bc-b\cdot tb}{a},\]
so
\[PQ+PE+QF=ta+\dfrac{2bc-t\left(b^2+c^2\right)}{a}=ta+\dfrac{2bc}{a}-ta=\boxed{\dfrac{2bc}{\sqrt{b^2+c^2}}},\]
independent of the position of $PQ$.
